\chapter{创新点与应用价值}
\chapterlead{我们没有发明新的模型，做的是把模型放进一套能被检查的流程里。}

\section{创新点}

\subsection{理解与计算彻底分开，并且由程序强制执行}
“让模型调用计算器”是常见做法，但多数系统只是建议模型这样做。\sysname{} 的不同在于强制：模型抄的数要逐字回原文找到，正文的数要按展示精度对回产物，
改写不能新增数字。模型想心算也没有用，算出来的数过不了护栏。

\subsection{双通道抽取，冲突即阻断}
用一条不依赖模型的规则通道独立解析表格，与模型抽取互相佐证。两条通道出错的方式不同，一致地错的概率很低；不一致时阻断，而不是替任何一方做选择。
在三家真实年报上，规则通道佐证了 93.0\% 的模型核验值，错配始终为 0。

\subsection{按展示精度定容差}
无论是数字护栏还是纠错核查，容差都由研报数字自身的写法决定。在 1,334 条注入断言上，这一策略同时做到了零误报和全召回，而任何固定容差都做不到。

\subsection{上下文隔离的多会话流水线与程序闸门}
五个角色各是一次独立会话，只靠文件交接，每一步由程序判定是否完成。闸门按状态而不是按会话退出码推进，异常中断后可以从断点继续。
评分标准与评测结果对生成方不可见，防止“应试”。

\subsection{会给自己打分、也会检查打分者的质检系统}
四层评分标准把能程序判的部分全部交给程序；模型评审用重复评审、人为缺陷和人工复核三组证据校准，并据此修订评分标准。
我们把评审做得不好的地方也写了出来（人为制造缺陷的判出率只有 42\%，整节删除类几乎判不出），而不是只报总分。

\section{应用价值}

\paragraph{对研究员。}一家新公司的年报，一个多小时可以得到一份数字全部可追溯的初稿、一个公式联动的估值模型和一份列明疑点的复核清单。
研究员可以把时间花在复核清单和估值假设上，这正是最需要人判断的地方。

\paragraph{对质控与合规。}纠错核查可以单独使用：输入一份研报和对应年报，输出逐条核查表，错误类型和修改建议具体到页码。
误报率低是它能被实际使用的前提。

\paragraph{对教学。}每个数字的来历、每个估值假设的依据、每一步计算都摊开在文件里，适合作为财务分析和估值课程的案例材料。
学生可以在工作簿里改一个假设，看三张表和每股价值怎么变。

\paragraph{可推广性。}计算引擎与具体模型无关，更换模型只需替换会话启动命令；技能说明不含行业模板，新公司不需要改代码。
目前的边界是：只覆盖非金融企业的 A 股年报与招股书；季报只走模型通道；扫描版 PDF 直接拒收，不做 OCR。
